%% 
%% Copyright 2007-2020 Elsevier Ltd
%% 
%% This file is part of the 'Elsarticle Bundle'.
%% ---------------------------------------------
%% 
%% Template article for Elsevier's document class `elsarticle`
%% with harvard style bibliographic references

\documentclass[final,5p,times,twocolumn,authoryear]{elsarticle}

%% The amssymb package provides various useful mathematical symbols
\usepackage{amssymb}
\usepackage{lipsum}
\usepackage{hyperref}

%% Comandos abreviados personalizados
\newcommand{\kms}{km\,s$^{-1}$}
\newcommand{\msun}{$M_\odot$}

\journal{Physics Open}

\begin{document}

% ==========================================
% PORTADA (Según Lineamientos UTEC)
% ==========================================
\begin{titlepage}
    \centering
    \vspace*{2cm} % Espacio superior
    
    {\huge\bfseries Efecto del método de integración numérica en modelos de flujo rectificado sobre distribuciones bidimensionales \par}
    
    \vspace{2.5cm}
    
    {\Large\textbf{Integrantes del grupo:}\par}
    \vspace{0.5cm}
    {\large
    Tommy Bryan Loo Leon - [202510241]  \par
    Jesús Valentín Niño Castañeda - [202310452]  \par
    Leonardo Jesús Medina Gago - [202410022] \par
    Lemuel Uziel Ramos Urbano - [202520133]  \par
    Fabricio Joaquin Rueda Zapata - [202510114]  \par
    }
    
    \vspace{2.5cm}
    
    {\Large\textbf{Curso:}\par}
    {\large CC2104 Métodos Numéricos \par}
    
    \vspace{2cm}
    
    {\Large\textbf{Profesor:}\par}
    {\large [Patricia Reynoso Quispe] \par}
    
    \vfill
    
    {\Large\textbf{Fecha de entrega:}\par}
    {\large 9 de Septiembre de 2026 \par}
    
    \vspace*{2cm}
\end{titlepage}
% ==========================================
% FIN DE LA PORTADA
% ==========================================

\begin{frontmatter}

\title{Efecto del método de integración numérica en modelos de flujo rectificado sobre distribuciones bidimensionales}

\author[first]{Tommy Bryan Loo Leon}
\author[first]{Jesús Valentín Niño Castañeda}
\author[first]{Leonardo Jesús Medina Gago}
\author[first]{Lemuel Uziel Rambos Urbano}
\author[first]{Fabricio Joaquin Rueda Zapata}

\affiliation[first]{organization={UTEC},%
            addressline={}, 
            city={},
            postcode={}, 
            state={},
            country={}}

\begin{abstract}
The project will study a continuous flow described by a differential equation that transports samples from a simple initial distribution to a two-dimensional target distribution. Instead of training a large-scale generative image model, we will implement different test vector fields or, alternatively, a small neural network to approximate the transport field and observe how the distribution evolves over time. We will use synthetic data and evaluate how the integrator affects the trajectory, final error, and computational cost. The ultimate goal will be to evaluate the relationship between accuracy and efficiency, identifying how important it is for the trajectories to be straight or close to the ideal trajectories. The theoretical foundation of the project will be related to papers that present different variants of rectified flow, flow matching, and straighter transport trajectories.
\end{abstract}

\begin{keyword}
Flow Matching \sep Rectified Flow \sep Integración Numérica \sep Modelos Generativos
\end{keyword}

\end{frontmatter}

%% main text

\section{Introduction}
\label{introduction}

Muchos problemas modernos requieren transformar un con-
junto de puntos, generados de manera simple y aleatoria, en
otro conjunto de puntos que sigue una distribución mucho más
compleja. Este tipo de transformación puede describirse me-
diante una ecuación diferencial ordinaria que indica, en cada
instante de tiempo, la velocidad con la que cada punto debe de-
splazarse para alcanzar su posición final. Dado que estas ecua-
ciones diferenciales rara vez admiten una solución analítica ex-
acta, resulta necesario resolverlas mediante métodos numéricos
de integración, cuya elección determina de manera directa tanto
la precisión de la trayectoria obtenida como el costo computa-
cional del proceso. Este problema resulta particularmente rele-
vante en el contexto de los modelos generativos modernos, en
los cuales el transporte eficiente y preciso de distribuciones de
probabilidad constituye un componente central.


Para profundizar en el tema, se han seleccionado trabajos recientes y relevantes que abordan distintos enfoques y avances en Flow Matching y métodos relacionados. Estos antecedentes permiten comprender la evolución de las ideas y los desafíos que se han presentado en el área.

Como primer antecedente, \citet{lipman2022flow} introdujeron Flow Matching, un marco que permite entrenar campos vectoriales para modelar el transporte continuo entre una distribución inicial y una distribución objetivo mediante ecuaciones diferenciales ordinarias. Lo interesante de su propuesta es que posibilita definir diferentes trayectorias de probabilidad y entrenar el campo de velocidad sin necesidad de simular previamente dichas trayectorias durante el entrenamiento. Además, los autores muestran que los caminos basados en transporte óptimo pueden generar muestras de forma más rápida y eficiente en comparación con otras alternativas, gracias al uso de métodos numéricos para resolver las ODE en la fase de muestreo, lo que proporciona la base conceptual y matemática sobre la que se apoyan los modelos de transporte analizados en este proyecto.

Posteriormente, \citet{liu2022flow} propusieron Rectified Flow, cuyo objetivo es aprender trayectorias de transporte lo más rectas posible entre las distribuciones inicial y final. La idea principal consiste en emplear un campo de velocidad que aproxime la trayectoria lineal que une los puntos correspondientes de ambas distribuciones. Esta característica resulta especialmente interesante desde una perspectiva numérica, ya que una trayectoria recta permite una discretización temporal mucho más gruesa sin perder precisión. Los autores demostraron que, mediante el procedimiento llamado reflow, es posible obtener trayectorias cada vez más rectas y lograr buenos resultados incluso utilizando un solo paso de Euler, lo que implica que se pueden reducir considerablemente los pasos necesarios para la integración y, en consecuencia, el tiempo requerido para generar muestras.

En esta misma línea, \citet{pooladian2023multisample} desarrollaron Multisample Flow Matching, una extensión que introduce acoplamientos no triviales entre muestras de las distribuciones de origen y destino. El objetivo de esta propuesta es aprovechar la estructura existente entre las muestras para construir campos vectoriales con trayectorias más rectas. Los autores demostraron que, además de reducir la varianza durante el entrenamiento, este enfoque permite generar muestras de buena calidad utilizando un menor número de evaluaciones del campo vectorial, lo que evidencia experimentalmente que la geometría de las trayectorias puede influir de manera directa en la eficiencia computacional del proceso de integración, y resalta la importancia de considerar este aspecto en el presente estudio.

De manera relacionada, \citet{tong2023improving}  exploraron el uso del transporte óptimo dentro de Flow Matching a través de su propuesta llamada Optimal Transport Conditional Flow Matching. En su estudio, los autores muestran que los acoplamientos obtenidos mediante transporte óptimo permiten construir flujos más simples y estables, además de facilitar una inferencia más rápida. Estos resultados refuerzan la idea de que la manera en que se establece la correspondencia entre las distribuciones inicial y objetivo influye directamente en la complejidad del campo vectorial y, por lo tanto, en la dificultad de su resolución numérica, lo que resulta especialmente relevante para el análisis realizado en este proyecto.

Posteriormente, \citet{yang2024consistency}  presentaron Consistency Flow Matching, una propuesta enfocada en construir trayectorias rectas mediante la aplicación de restricciones de consistencia en el campo de velocidad. A diferencia de otros enfoques que dependen principalmente de procesos iterativos para rectificar las trayectorias, esta metodología incorpora de manera explícita la consistencia de la velocidad para enlazar distintos momentos en el tiempo con un mismo destino final. Los autores argumentan que este mecanismo favorece la simplicidad de las trayectorias y posibilita una generación eficiente, requiriendo menos evaluaciones del campo vectorial. Así, su trabajo subraya que mantener trayectorias menos complejas sigue siendo clave para optimizar la velocidad de inferencia en los modelos basados en Flow Matching, aportando un enfoque novedoso en este ámbito.

Desde el enfoque teórico, \citet{zhou2025error}  analizaron el error presente en los modelos generativos desarrollados con Flow Matching. Su trabajo se centra en cómo la distribución generada se aproxima a la distribución objetivo y en qué condiciones se puede garantizar esta convergencia, utilizando como referencia la distancia de Wasserstein-2. Lo interesante de este antecedente es que permite separar claramente el error que proviene del aprendizaje del campo de velocidad de aquel que surge al resolver numéricamente la ecuación diferencial. Esta distinción resulta crucial para este proyecto, pues ayuda a entender mejor cómo el método de integración influye en la distribución final que se obtiene.

Finalmente, \citet{xiao2026solvers}  llevó a cabo un análisis comparativo de distintos métodos de integración numérica aplicados a modelos de Flow Matching. En su estudio, exploró alternativas como Euler, midpoint explícito, Runge-Kutta de cuarto orden y Dormand-Prince, realizando pruebas con distribuciones sintéticas bidimensionales. Al comparar la calidad de las muestras y la cantidad de evaluaciones requeridas del campo vectorial, notó que los métodos más avanzados pueden ser más eficientes sin comprometer la calidad de los resultados. De esta manera, su trabajo muestra cómo la elección del método de integración puede marcar diferencias sutiles pero importantes en el comportamiento práctico de estos modelos.

En conjunto, estos antecedentes sugieren que la eficiencia de los modelos basados en Flow Matching y Rectified Flow no depende únicamente de la capacidad del campo vectorial para transformar una distribución en otra. También influyen aspectos como la geometría de las trayectorias y el método numérico empleado para resolver la ecuación diferencial. Por esto, resulta interesante analizar cómo se comporta el modelo al mantener constante el campo vectorial y variar únicamente el método de integración, lo que permite observar con mayor claridad cómo la discretización numérica afecta tanto la trayectoria como el resultado final y los recursos computacionales requeridos.


A pesar de los recientes y notables avances en los modelos generativos continuos, un inconveniente importante tanto de los Flujos Normalizantes Continuos (CNFs, descritos por EDOs) como de los modelos de difusión, tambien descritos por EDOs, es que su simulación e integración numérica requieren múltiples evaluaciones de la red neuronal para generar una muestra de alta calidad. Este proceso iterativo da como resultado tiempos de inferencia prolongados, lo que limita su eficiencia y aplicabilidad en escenarios de tiempo real. Para mitigar este inconveniente, investigaciones recientes han motivado el desarrollo de métodos que imponen propiedades de transporte óptimo (OT) en las EDOs neuronales, produciendo flujos con trayectorias más rectas que pueden integrarse con precisión utilizando un menor número de evaluaciones de la red neuronal. Sin embargo, la literatura carece de un análisis exhaustivo sobre cómo la elección específica del método de integración numérica altera la trayectoria, el error y el costo cuando se aplican estos flujos enderezados. El presente trabajo se justifica en la necesidad de llenar este vacío analizando el efecto del integrador numérico sobre modelos de flujo rectificado en un entorno controlado de distribuciones bidimensionales. Comprender la relación exacta entre la rectitud de la trayectoria, la precisión del integrador numérico y la eficiencia computacional es un paso indispensable para optimizar el despliegue de modelos generativos modernos.


Por lo expuesto, como objetivo general tenemos: modelar numéricamente el transporte continuo de muestras mediante ecuaciones diferenciales que transformen una distribución inicial simple en una distribución objetivo bidimensional, con el propósito de analizar la relación entre la precisión, la eficiencia computa cional y la rectitud de las trayectorias generadas mediante diferentes campos vectoriales y métodos de integración. En base a ello desarrollamos nuestros objetivos específicos los cuales son:
\begin{enumerate}
\item Formular e implementar un modelo numérico basado en EDOs que represente el transporte de muestras desde una distribución inicial simple hacia una distribución objetivo bidimensional mediante campos vectoriales de prueba.
\item Evaluar el comportamiento del modelo bajo un conjunto base de parámetros, considerando la trayectoria de las muestras, el error respecto a la distribución objetivo y el costo computacional asociado al proceso de integración.

\item Analizar la robustez del modelo mediante la variación de parámetros relevantes, comparando los cambios en la precisión, la rectitud de las trayectorias y el costo computacional obtenidos bajo diferentes configuraciones.

\item Aplicar el modelo numérico al transporte de distribuciones bidimensionales para determinar la relación entre la rectitud de las trayectorias, el método de integración empleado y la eficiencia en la generación de muestras.
\end{enumerate}

\section{Metodología}
\label{sec:metodologia}

\subsection{Modelo del Sistema}
El desarrollo de este proyecto se fundamenta en la formulación de Modelos Generativos Continuos (como Continuous Normalizing Flows), los cuales modelan la transformación progresiva de una densidad de probabilidad empírica a través del tiempo. Teóricamente, la evolución de esta densidad de probabilidad $p_t(x)$ en el tiempo $t \in [0,1]$ está gobernada rigurosamente por la ecuación de continuidad (o ecuación de transporte):
\begin{equation}
    \frac{\partial p_t(x)}{\partial t} + \nabla \cdot (p_t(x)v_t(x)) = 0
\end{equation}
donde el primer término representa la variación temporal de la densidad de probabilidad local, y el segundo término es la divergencia espacial del flujo de probabilidad. Esta Ecuación en Derivadas Parciales (EDP) asegura la conservación de la masa (probabilidad), determinando cómo el campo de velocidades $v_t(x)$ acumula o dispersa la distribución en el espacio.

Para evitar la alta complejidad computacional de resolver la EDP directamente, el sistema se aborda desde una perspectiva lagrangiana. El transporte continuo se describe de manera equivalente mediante el seguimiento de trayectorias individuales a través de la siguiente Ecuación Diferencial Ordinaria (EDO):
\begin{equation}
    \frac{dZ_t}{dt} = v_\theta(Z_t, t), \quad Z_0 \sim \pi_0
\end{equation}
El problema central consiste en encontrar un campo vectorial aproximado por una red neuronal, $v_\theta(Z_t, t)$, que permita transportar muestras desde una distribución inicial simple y tratable $\pi_0$ (ruido gaussiano) hacia la distribución objetivo $\pi_1$ (los datos a generar). La solución en el instante $t=1$, denotada como $Z_1$, sigue la distribución empírica de los datos.

Para entrenar $v_\theta$ sin resolver iterativamente la EDO durante la optimización, se emplean enfoques de emparejamiento libre de simulación (\textit{simulation-free}). Bajo el marco de Rectified Flow, el objetivo es conectar las muestras $X_0 \sim \pi_0$ y $X_1 \sim \pi_1$ mediante una interpolación lineal $X_t = tX_1 + (1-t)X_0$. La red se optimiza resolviendo el problema de mínimos cuadrados:
\begin{equation}
    \min_\theta \int_0^1 \mathbb{E} \left[ || (X_1 - X_0) - v_\theta(X_t, t) ||^2 \right] dt
\end{equation}
Equivalentemente, el marco de Flow Matching define una ruta de probabilidad condicional basada en Transporte Óptimo. Ambos enfoques garantizan matemáticamente que el transporte de partículas ocurra en líneas rectas y a velocidad constante, lo que simplifica la topología del campo vectorial y reduce drásticamente la dificultad de la integración numérica durante la fase de inferencia.

\subsection{Supuestos y Simplificaciones}
Para mantener un entorno controlado que permita enfocar el análisis puramente en el desempeño numérico, se establecen las siguientes simplificaciones y supuestos:
\begin{itemize}
    \item \textbf{Dimensionalidad reducida:} Se asume que el análisis numérico sobre distribuciones empíricas bidimensionales (conjuntos de datos sintéticos en 2D, como \textit{moons} o \textit{8-gaussians}) es representativo de la dinámica subyacente. Esto simplifica la visualización de las trayectorias de las EDOs y reduce significativamente la carga computacional, permitiendo evaluar múltiples métodos numéricos de forma eficiente.
    \item \textbf{Distribución base estándar:} Se asume como distribución inicial tratable $\pi_0$ a una distribución Normal multivariada estándar $\mathcal{N}(0, I)$, facilitando el muestreo inicial $Z_0$.
    \item \textbf{Precisión del Campo Vectorial:} Se asume que la red neuronal implementada (un Perceptrón Multicapa) aproxima con suficiente exactitud el campo vectorial ideal $v(x, t)$. En consecuencia, cualquier error en la distribución final generada $Z_1$ se atribuirá predominantemente al error de discretización del integrador numérico y no a un déficit en el entrenamiento del modelo.
    \item \textbf{Ventaja de las trayectorias rectas:} Se asume teóricamente que la rectitud de los flujos inducidos por Flow Matching y Rectified Flow reduce significativamente la curvatura del campo vectorial. Bajo este supuesto, el error de truncamiento local de los métodos numéricos será menor, permitiendo el uso de tamaños de paso ($dt$) más grandes sin degradar drásticamente la calidad de la solución.
\end{itemize}

\subsection{Métodos Numéricos Propuestos}
Una vez que el modelo ha sido entrenado y el campo vectorial $v_\theta(Z_t, t)$ se encuentra definido, la fase de generación (inferencia) requiere aproximar la solución de la EDO desde $t=0$ hasta $t=1$. Para este propósito, se propone implementar y evaluar comparativamente tres métodos numéricos de un paso para problemas de valor inicial (PVI).

Acorde a los lineamientos del curso, se proponen \textbf{dos métodos numéricos revisados en clase}:
\begin{enumerate}
    \item \textbf{Método de Euler (Primer Orden):} Se implementa como el método base de comparación (baseline). Al ser el método explícito más sencillo y de paso fijo, se espera que presente un alto error de truncamiento frente a tamaños de paso grandes, permitiendo evaluar qué tan permisivas son las trayectorias rectas frente a integradores de bajo orden computacional.
    \item \textbf{Método de Runge-Kutta de Cuarto Orden (RK4):} Se implementa como el integrador de alta precisión revisado en el curso. Su evaluación permitirá establecer un estándar de alta calidad visual de las muestras generadas de forma estática. Dado su mayor costo por iteración (cuatro evaluaciones de la función $v_\theta$ por paso temporal fijo), servirá para analizar críticamente el compromiso entre costo computacional (Número de Evaluaciones de la Función, NFE) y precisión.
\end{enumerate}

Para complementar el análisis, se propone implementar \textbf{un método numérico no revisado explícitamente en clase}:
\begin{enumerate}
    \setcounter{enumi}{2}
    \item \textbf{Método de Runge-Kutta Adaptativo (RK45 - Fehlberg / Dormand-Prince):} A diferencia de los métodos de paso fijo, este método empotrado calcula simultáneamente una aproximación de cuarto y quinto orden para estimar el error de truncamiento local en cada iteración. Esto permite adaptar dinámicamente el tamaño del paso temporal $dt$. Se propone este método porque es altamente eficiente para EDOs bien comportadas; en las zonas donde las trayectorias del modelo generativo sean efectivamente rectas, el RK45 aumentará automáticamente el tamaño de paso, optimizando el costo computacional (minimizando el NFE) sin comprometer la estricta tolerancia de error ni la precisión de la distribución generada $Z_1$.
\end{enumerate}


% \section{Title 2}
% \lipsum[1]
% 
% \subsection{Subsection title}
% 
% \begin{figure}[htbp]
% 	\centering     \includegraphics[width=0.4\textwidth, angle=-90]{PHYSO_cover_image.pdf}	
% 	\caption{Physics Open journal cover} 
% 	\label{fig_mom0}%
% \end{figure}
% 
% A random equation, the Toomre stability criterion:
% 
% \begin{equation}
%     Q = \frac{\sigma_v \times \kappa}{\pi \times G \times \Sigma}
% \end{equation}
% 
% \section{Title 3}
% \lipsum[2]
% 
% \subsection{Subsection title}
% \lipsum[3]
% 
% \begin{table}[htbp]
% \centering
% \begin{tabular}{l c c c} 
%  \hline
%  Source & RA (J2000) & DEC (J2000) & $V_{\rm sys}$ \\
%          & [h,m,s]    & [o,','']    & \kms          \\
%  \hline
%  NGC\,253 & 00:47:33.120 & -25:17:17.59 & $235 \pm 1$ \\ 
%  M\,82 & 09:55:52.725, & +69:40:45.78 & $269 \pm 2$  \\ 
%  \hline
% \end{tabular}
% \caption{Random table with galaxies coordinates and velocities.}
% \label{Table1}
% \end{table}
% 
% 
% \section{Discussion}
% \lipsum[4]
% 
% \section{Summary and conclusions}
% \lipsum[1-4]
% 
% 
% \section*{Acknowledgements}
% Thanks to ...
% 
% \appendix
% 
% \section{Appendix title 1}
% 
% \section{Appendix title 2}

%% Bibliografía integrada directamente
\begin{thebibliography}{00}

\bibitem[Lipman et al.(2022)]{lipman2022flow}
Lipman, Y., Chen, R.~T.~Q., Ben-Hamu, H., Nickel, M., Le, M., 2022.
Flow matching for generative modeling.
arXiv preprint arXiv:2210.02747.
\url{https://doi.org/10.48550/arXiv.2210.02747}

\bibitem[Liu et al.(2022)]{liu2022flow}
Liu, X., Gong, C., Liu, Q., 2022.
Flow straight and fast: Learning to generate and transfer data with rectified flow.
arXiv preprint arXiv:2209.03003.
\url{https://doi.org/10.48550/arXiv.2209.03003}

\bibitem[Pooladian et al.(2023)]{pooladian2023multisample}
Pooladian, A.-A., Ben-Hamu, H., Domingo-Enrich, C., Amos, B., Lipman, Y., Chen, R.~T.~Q., 2023.
Multisample flow matching: Straightening flows with minibatch couplings.
In: Proceedings of the 40th International Conference on Machine Learning. PMLR, pp. 28100--28120.
\url{https://proceedings.mlr.press/v202/pooladian23a.html}

\bibitem[Tong et al.(2023)]{tong2023improving}
Tong, A., Fatras, K., Malkin, N., Huguet, G., Zhang, Y., Rector-Brooks, J., Wolf, G., Bengio, Y., 2023.
Improving and generalizing flow-based generative models with minibatch optimal transport.
arXiv preprint arXiv:2302.00482.
\url{https://doi.org/10.48550/arXiv.2302.00482}

\bibitem[Yang et al.(2024)]{yang2024consistency}
Yang, L., Zhang, Z., Zhang, Z., Liu, X., Xu, M., Zhang, W., Meng, C., Ermon, S., Cui, B., 2024.
Consistency flow matching: Defining straight flows with velocity consistency.
arXiv preprint arXiv:2407.02398.
\url{https://doi.org/10.48550/arXiv.2407.02398}

\bibitem[Zhou and Liu(2025)]{zhou2025error}
Zhou, Z., Liu, W., 2025.
An Error Analysis of Flow Matching for Deep Generative Modeling.
In: Proceedings of the 42nd International Conference on Machine Learning.
PMLR 267, 78903--78932.
\url{https://proceedings.mlr.press/v267/zhou25l.html}

\bibitem[Xiao(2026)]{xiao2026solvers}
Xiao, H., 2026.
From Euler to Dormand-Prince: ODE Solvers for Flow Matching Generative Models.
arXiv preprint arXiv:2605.00836.
\url{https://arxiv.org/abs/2605.00836}

\end{thebibliography}

\end{document}
